% =============================================================================
%  2.2 Хуки (Hooks)
% =============================================================================
\section{Хуки (Hooks)}

\term{Хуки}{Hooks} — механизм, позволяющий функциональным компонентам
использовать состояние и другие возможности React без написания классов.
Хуки были введены в React~16.8 и стали основным способом работы
с состоянием в современных приложениях.

Внутри React хуки хранятся как связный список, привязанный к
файбер-узлу компонента.  Порядок вызова хуков должен быть одинаковым
при каждом рендере — отсюда правило: хуки нельзя вызывать условно.

\subsection{useState}

\begin{code}{javascript}
function Counter() {
  const [count, setCount] = useState(0);

  return (
    <button onClick={() => setCount(c => c + 1)}>
      Счётчик: {count}
    </button>
  );
}
\end{code}

\subsection{useEffect}

\subsubsection{Массив зависимостей}

Число повторных вызовов эффекта зависит от массива зависимостей.
Если зависимость — примитив, сравнение выполняется за $O(1)$;
для массива из $d$ зависимостей стоимость сравнения:
$C = \sum_{i=1}^{d} O(1) = O(d)$.

\begin{equation}
  \text{Re-runs}(e) = \bigl|\{r \in \text{renders} : \text{deps}(e, r) \neq \text{deps}(e, r-1)\}\bigr|
\end{equation}

\begin{code}{javascript}
useEffect(() => {
  const controller = new AbortController();
  fetch(`/api/user/${id}`, { signal: controller.signal })
    .then(r => r.json())
    .then(setUser);
  return () => controller.abort();   // cleanup
}, [id]);
\end{code}

\note{Функция очистки (\inl{cleanup}) вызывается перед каждым
повторным запуском эффекта и при размонтировании компонента.
Используйте её для отмены запросов, отписки от событий и
освобождения ресурсов.}

\important{Пропуск зависимости в массиве \inl{useEffect} — одна из
самых частых ошибок в React-приложениях.  Используйте линтер-плагин
\inl{eslint-plugin-react-hooks} для автоматической проверки.}
